\begin{table*}[t]\centering\footnotesize
\caption{Research questions, systematization findings and evidence. Experimental follow-ups test the comparison rules; the SoK contribution is the record-level synthesis.}\label{tab:rq-findings}
\begin{tabularx}{\textwidth}{@{}p{.20\textwidth}YY@{}}\toprule
Question & Finding & Evidence in the main paper\\\midrule
RQ1: targets and observables & A common signal can serve confidentiality, attribution, evasion or authorized diagnostics; polarity changes the protected interest. & Complete core-family map (\S\ref{sec:corpus-map}, Table~\ref{tab:core-family-map}).\\
RQ2: origin versus access & Generating layers do not determine receiver access. Upstream protection need not survive a plaintext-aware downstream emitter; read cover does not imply safe writes. & Record lineage and three literature cases (\S\ref{sec:scope}--\ref{sec:discussion}); E5--E7.\\
RQ3: portability & A task holdout does not establish topology, hardware or model transfer; mixed configurations can preserve shortcuts. & Generator-specific holdouts (\S\ref{sec:evaluation}); customer-grouped analysis (\S\ref{sec:r5}).\\
RQ4: protection, state and cost & Evaluate leakage, logical state preservation and resource cost together. Full LLM utility under timing-dependent replanning is not established. & Twenty-family audit; E9--E14 leakage/state/cost contrasts (\S\ref{sec:controlled-placement}).\\\bottomrule
\end{tabularx}\end{table*}
